\section{RESULTS}

\subsection{Dataset Characteristics}

The preprocessed dataset contained 1,452 student records from module BBB, presentation 2013J, of which 959 were Not At Risk (66.05\%) and 493 were At Risk (33.95\%). This roughly two-to-one class ratio was preserved across the training and test sets through stratified splitting.

\subsection{Expert System Knowledge Base}

The full knowledge base is shown in Table~\ref{tab:knowledge-base}. The support column reports the number of training samples that reached each decision-tree leaf before the rule was converted into an expert-system rule.

\begin{small}
\begin{longtable}{p{0.65cm}p{6.9cm}p{2.2cm}p{1.1cm}p{1.2cm}}
\caption{Generated Expert System Knowledge Base}\label{tab:knowledge-base}\\
\toprule
Rule & Conditions & Prediction & Conf. & Support \\
\midrule
\endfirsthead
\toprule
Rule & Conditions & Prediction & Conf. & Support \\
\midrule
\endhead
\bottomrule
\endfoot
1 & \texttt{average\_CA\_score} $\leq$ 24.628 & At Risk & 1.00 & 111 \\
2 & 24.628 $<$ \texttt{average\_CA\_score} $\leq$ 29.174 & At Risk & 0.91 & 13 \\
3 & 29.174 $<$ \texttt{average\_CA\_score} $\leq$ 33.395 & At Risk & 1.00 & 13 \\
4 & 33.395 $<$ \texttt{average\_CA\_score} $\leq$ 44.616 AND \texttt{total\_clicks} $\leq$ 54.5 & At Risk & 1.00 & 16 \\
5 & 33.395 $<$ \texttt{average\_CA\_score} $\leq$ 44.616 AND \texttt{total\_clicks} $>$ 54.5 & At Risk & 0.86 & 90 \\
6 & 44.616 $<$ \texttt{average\_CA\_score} $\leq$ 54.105 AND \texttt{date\_registration} $\leq$ -20.5 & At Risk & 0.76 & 75 \\
7 & 44.616 $<$ \texttt{average\_CA\_score} $\leq$ 54.105 AND \texttt{date\_registration} $>$ -20.5 & Not At Risk & 1.00 & 6 \\
8 & 54.105 $<$ \texttt{average\_CA\_score} $\leq$ 61.988 AND \texttt{imd\_band} is not 90--100\% & At Risk & 0.56 & 87 \\
9 & 61.988 $<$ \texttt{average\_CA\_score} $\leq$ 75.105 AND \texttt{imd\_band} is not 90--100\% & Not At Risk & 0.66 & 348 \\
10 & 54.105 $<$ \texttt{average\_CA\_score} $\leq$ 75.105 AND \texttt{imd\_band} is 90--100\% & Not At Risk & 1.00 & 23 \\
11 & 75.105 $<$ \texttt{average\_CA\_score} $\leq$ 75.860 & Not At Risk & 1.00 & 24 \\
12 & 75.860 $<$ \texttt{average\_CA\_score} $\leq$ 80.663 & Not At Risk & 0.81 & 162 \\
13 & \texttt{average\_CA\_score} $>$ 80.663 AND \texttt{imd\_band} is not 10--20\% & Not At Risk & 0.97 & 179 \\
14 & \texttt{average\_CA\_score} $>$ 80.663 AND \texttt{imd\_band} is 10--20\% & Not At Risk & 0.76 & 14 \\
\end{longtable}
\end{small}


As shown in the decision tree, \texttt{average\_CA\_score} was the main factor used to separate At Risk and Not At Risk students. Very low scores were classified as At Risk with high confidence, while high scores were generally classified as Not At Risk. The middle score ranges were less certain and required additional conditions, such as \texttt{total\_clicks}, \texttt{date\_registration}, and \texttt{imd\_band}, showing that engagement, registration timing, and background factors were mainly used to refine predictions when assessment performance alone did not clearly indicate risk.

\subsection{Expert System Performance Against the ML Baselines}

Table~\ref{tab:model-performance} shows how these results compare to the two ML baselines.

\begin{table}[H]
\centering
\caption{Test Set Performance --- Positive Class: At Risk}
\label{tab:model-performance}
\begin{tabular}{lcccc}
\toprule
System & Accuracy & Precision & Recall & F1-score \\
\midrule
Expert System (ES) & 0.8316 & 0.7404 & 0.7778 & 0.7586 \\
Logistic Regression (LR) & 0.8282 & 0.7379 & 0.7677 & 0.7525 \\
Support Vector Machine (SVM) & 0.8488 & 0.7835 & 0.7677 & 0.7755 \\
\bottomrule
\end{tabular}
\end{table}


Table~\ref{tab:confusion-matrices} shows the corresponding confusion-matrix counts for the same held-out test set.

\begin{table}[H]
\centering
\caption{Confusion-Matrix Counts for the Held-Out Test Set}
\label{tab:confusion-matrices}
\begin{tabular}{llrr}
\toprule
System & Actual Class & Predicted Not At Risk & Predicted At Risk \\
\midrule
Expert System (ES) & Not At Risk & 165 & 27 \\
Expert System (ES) & At Risk & 22 & 77 \\
Logistic Regression (LR) & Not At Risk & 165 & 27 \\
Logistic Regression (LR) & At Risk & 23 & 76 \\
Support Vector Machine (SVM) & Not At Risk & 171 & 21 \\
Support Vector Machine (SVM) & At Risk & 23 & 76 \\
\bottomrule
\end{tabular}
\end{table}


The expert system achieved the highest recall, correctly identifying 77 of the 99 at-risk students in the test set. Logistic regression identified 76 at-risk students, and the support vector machine also identified 76. The support vector machine achieved the highest accuracy, precision, and F1-score, indicating a slightly stronger balance between false positives and false negatives. The differences between the three systems were small, so the results should be interpreted as metric-specific rather than as one system clearly outperforming the others on all criteria. The main advantage of the expert system is that it reached comparable performance while keeping the decision process available as explicit IF-THEN rules.

\subsection{Comparison with Related OULAD At-Risk Models}

Table~\ref{tab:related-work-comparison} compares our ES/LR/SVM results against the at-risk results reported by \textcite{ujkani2024oulad}.

\begin{table}[H]
\centering
\caption{Contextual Comparison with OULAD At-Risk Results Reported in Related Work}
\label{tab:related-work-comparison}
\resizebox{\textwidth}{!}{%
\begin{tabular}{llcccc}
\toprule
Study & Model & Accuracy (\%) & Precision$_{risk}$ & Recall$_{risk}$ & F1$_{risk}$ \\
\midrule
Our work & Expert System (ES) & 83.2 & 0.740 & 0.778 & 0.759 \\
Our work & Logistic Regression (LR) & 82.8 & 0.738 & 0.768 & 0.753 \\
Our work & SVM & 84.9 & 0.784 & 0.768 & 0.776 \\
\midrule
\cite{ujkani2024oulad} & Random Forest & 86.0 & 0.860 & 0.970 & 0.910 \\
\cite{ujkani2024oulad} & Gradient Boosting & 81.0 & 0.830 & 0.960 & 0.890 \\
\cite{ujkani2024oulad} & k-NN & 82.0 & 0.840 & 0.940 & 0.890 \\
\cite{ujkani2024oulad} & MLP & 85.0 & 0.880 & 0.940 & 0.910 \\
\cite{ujkani2024oulad} & CNN & 86.0 & 0.870 & 0.950 & 0.910 \\
\cite{ujkani2024oulad} & LSTM & 81.0 & 0.840 & 0.920 & 0.880 \\
\cite{ujkani2024oulad} & XGBoost & 86.0 & 0.870 & 0.960 & 0.910 \\
\cite{ujkani2024oulad} & Custom NN & 93.0 & 0.950 & 0.970 & 0.960 \\
\bottomrule
\end{tabular}%
}
\vspace{1mm}
\footnotesize{\textit{Note:} Related-work metrics are taken from Tables 7--10 of \textcite{ujkani2024oulad}. Their at-risk task reports Pass and Fail, where the Fail class includes non-success outcomes used for identifying at-risk students. The comparison is contextual because the studies differ in scope, feature setup, and modelling aim.}
\end{table}


This comparison should be read as contextual rather than directly equivalent. Ujkani et al. used OULAD for a broader modelling study and reported results for several machine learning and deep learning models, while this project used only BBB/2013J and focused on a transparent rule-based expert system. Their reported models generally achieved stronger risk-class metrics, especially the custom neural network. The expert system in this project therefore should not be presented as outperforming those models. Its contribution is narrower: within the BBB/2013J early-window setup, it achieved performance close to the two local ML baselines while producing a compact, inspectable rule base.
